\chapter{Introduction}
\label{ch:intro}

\section{Problem Statement}
\label{sec:intro-problem}

\paragraph{Context.} Automated program repair (APR) fixes a faulty program
without a human writing the patch~\citep{legoues2019apr}. Most tools work in a
loop: they generate a candidate patch and run the test suite on it; a candidate
that passes every test is called \emph{plausible}~\citep{qi2015plausibility}.
Today the candidate comes from a large language model (LLM). If a test fails,
the tool shows the failure to the model and asks again~\citep{xia2024chatrepair}.
Each round of this loop has two costs (Figure~\ref{fig:intro-bottleneck}): the
model call and the test run.

\paragraph{Need.} The cost per bug is now reported next to the repair rate:
\textsc{ChatRepair} spends about \$0.42 per bug~\citep{xia2024chatrepair} and
\textsc{RepairAgent} about $270{,}000$ tokens per bug, with $99\,\%$ of its time
spent running tools, mostly tests~\citep{bouzenia2025repairagent}. Tools already
try to lower this cost~\citep{hidvegi2024cigar}, and long, costly sessions are
named as an obstacle to using LLM-based repair in
practice~\citep{yang2025llmrepairsurvey}. The practical need is therefore to \emph{lower the cost per repaired bug without
repairing fewer bugs}. A typical setting is a repair bot in continuous
integration (CI), which is where developers most want repair to
run~\citep{noller2022trust}. Such a bot must respond within a bounded time and
budget, fit into the existing pipeline, never propose a patch that the tests
reject, and explain its failures.

\begin{figure}[htbp]
\centering
\definecolor{okgreen}{HTML}{2E7D32}
\resizebox{0.97\textwidth}{!}{%
\begin{tikzpicture}[
  font=\small,
  box/.style={draw=black!55,rounded corners=2pt,align=center,inner sep=5pt,
              text width=38mm,minimum height=16mm,fill=black!3},
  res/.style={draw=okgreen,line width=0.8pt,rounded corners=2pt,align=center,
              inner sep=4pt,text width=20mm,minimum height=11mm,fill=okgreen!10},
  cost/.style={font=\scriptsize\bfseries,text=bkblue,align=center},
  tp/.style={draw=okgreen,fill=okgreen!45,minimum width=3.6mm,minimum height=3.6mm,
             inner sep=0pt,rounded corners=0.5pt},
  tf/.style={tp,draw=pendred,fill=pendred!60},
  tn/.style={tp,draw=black!30,fill=white},
  waste/.style={circle,draw=pendred,fill=pendred!8,text=pendred,font=\scriptsize\bfseries,
                inner sep=1pt,minimum size=4.6mm},
  arr/.style={-{Stealth[length=2.4mm]},draw=black!60}
]
\node (llm) at (0,0)    [box] {\textbf{LLM}\\[2pt]{\scriptsize proposes a candidate patch}};
\node (val) at (7.2,0)  [box] {\textbf{Validation}\\[2pt]{\scriptsize builds it and runs the tests,}\\{\scriptsize stopping at the first failure}};
\node (pl)  at (12.6,0) [res] {plausible\\ patch};
\node[cost,above=2pt of llm] {cost 1: model call\\ (tokens, latency)};
\node[cost,above=2pt of val] {cost 2: validation\\ (CPU time, wall-clock)};
\draw[arr] (llm.east) -- node[below,font=\scriptsize]{candidate} (val.west);
\draw[arr,draw=okgreen] (val.east) -- node[above,font=\scriptsize,text=okgreen]{all pass} (pl.west);
\coordinate (fb) at ([xshift=-14mm]val.south);
\coordinate (fbl) at ([yshift=-15mm]fb);
\draw[arr] (fb) -- (fbl) -| node[pos=0.25,below,font=\scriptsize]
  {a test fails: next round, the failure is fed back} (llm.south);
\begin{scope}[shift={([xshift=-6mm,yshift=-6mm]val.south)}]
  \foreach \i in {0,...,3} { \node[tp] at ({0.5*\i},0) {}; }
  \node[tf] at (2.0,0) {};
  \foreach \i in {5,6} { \node[tn] at ({0.5*\i},0) {}; }
  \node[tp,minimum width=2.8mm,minimum height=2.8mm] (lp) at (0,-0.6) {};
  \node[font=\scriptsize,anchor=west] at (lp.east) {pass};
  \node[tf,minimum width=2.8mm,minimum height=2.8mm] (lf) at (1.1,-0.6) {};
  \node[font=\scriptsize,anchor=west] at (lf.east) {fail};
  \node[tn,minimum width=2.8mm,minimum height=2.8mm] (ln) at (2.1,-0.6) {};
  \node[font=\scriptsize,anchor=west] at (ln.east) {not run};
  \node[waste] (w2) at (3.6,0) {2};
\end{scope}
\node[waste] (w1) at ($(llm.east)!0.5!(val.west)+(0,0.34)$) {1};
\node[waste] (w3) at ($(llm.south |- fbl)!0.5!(fbl)+(0,0.34)$) {3};
\end{tikzpicture}}
\caption[One round of an LLM-based repair loop]{One round of an LLM-based repair
loop and its two costs. Under the validation step, the tests of a rejected
candidate: passing tests (green) run before the failing one (red), and the rest
are not run. The marks are the three wastes: \textcircled{\scriptsize 1} a
repeated candidate is generated and tested again; \textcircled{\scriptsize 2}
passing tests run before the one that rejects a wrong candidate;
\textcircled{\scriptsize 3} rounds go on after the model has stopped making
progress.}
\label{fig:intro-bottleneck}
\end{figure}

\paragraph{Problem.} A large part of this cost is spent without benefit, in
three ways:
\begin{itemize}[leftmargin=*,itemsep=2pt]
\item \textbf{Repeats.} The model proposes patches it has already proposed:
$65\,\%$ of the patches of \textsc{ChatRepair} were duplicates~\citep{chen2025conversational}.
Each repeat costs a model call and often a test run.
\item \textbf{Slow rejections.} One failing test is enough to reject a wrong
patch, but a fixed test order may run many passing tests first. Test orderings
for repair exist~\citep{qi2013frtp}, but they ignore how long each test takes.
\item \textbf{Stuck sessions.} A session goes on after the model has stopped
making progress, until its budget runs out. These failed sessions are the
expensive ones~\citep{fan2025sweeffi}.
\end{itemize}
The size of each waste is unknown, because existing studies report either the
generation or the validation cost, for different tools and settings. It is
therefore not clear which part of the cost is worth reducing.

\paragraph{Goal.} This thesis first measures how the cost of a repair loop is distributed. It
then builds a \emph{validation layer}, a component placed between the repair
loop and the test runner, that addresses the three wastes:
\begin{enumerate}[leftmargin=*,itemsep=2pt]
\item \textbf{Module M1, Repeat Guard}, avoids validating the same patch twice: it recognises a repeated patch, reuses
its result and informs the model;
\item \textbf{Module M2, Test Prioritiser}, rejects wrong patches early: it runs first the tests with the highest
probability of failing per second of running time;
\item \textbf{Module M3, Budget Guard}, stops unproductive sessions: it ends a session when the model makes no progress or the budget is spent, and reports the reason.
\end{enumerate}
The first two never change whether a patch is accepted. The third can, so its
effect on the repair rate is measured.

\section{Research Questions and Contributions}
\label{sec:intro-this}

\begin{itemize}[leftmargin=*,itemsep=3pt]
\item \textbf{RQ1: Where does the cost of an LLM-based repair session go?}
The cost per repaired bug, split into generating and validating patches, and
how much of it falls after the model is stuck.
\item \textbf{RQ2: How much does each module save?} Each of the three
modules is measured on the part of the cost it targets. The test ordering,
\km, is compared with \frtp~\citep{qi2013frtp} and with the coverage-based rule
of the mutation tool PIT~\citep{coles2016pit}.
\item \textbf{RQ3: Does the whole layer lower the cost per repaired bug?}
Against a standard conversational loop, without a lower repair rate.
\item \textbf{RQ4: Why does \km\ differ from existing orderings?} Which of
its parts matters, and when.
\end{itemize}

\paragraph{Contributions.}
\begin{enumerate}[label=\textbf{C\arabic*},leftmargin=*,itemsep=2pt]
\item \textbf{A cost breakdown of LLM-based repair loops}, with the most each
kind of technique could save.
\item \textbf{A cost-aware validation layer} of three modules: a repeat guard,
a test prioritiser and a budget guard.
\item \textbf{\km} (Adaptive Rate-over-Cost), the test ordering at the core of
the prioritiser: it
learns, during the session, how likely each test is to fail and how long it
takes, and needs no extra instrumentation.
\item \textbf{A two-level evaluation} on Defects4J: a simulation in which every
ordering sees the same patches, and live runs with a real model.
\end{enumerate}
C1, C3 and C4 are the main research contributions. The repeat guard and the
budget guard are included mainly for practical reasons, since they are simple,
do not affect correctness and are absent from current tools (Section~\ref{sec:rel-practice}).

\section{Scope}
\label{sec:intro-scope}

\begin{table}[htbp]
\centering
\caption{Scope of the study.}
\label{tab:intro-scope}
\footnotesize
\begin{tabularx}{\textwidth}{@{}>{\raggedright\arraybackslash}p{0.2\textwidth} >{\raggedright\arraybackslash}X@{}}
\toprule
In scope & Defects4J~\citep{just2014defects4j} (Java), $100$ single-file bugs; one LLM through an API; a conversational loop with test feedback; the faulty methods are given to the model; success means a plausible patch\\
Measured & time, CPU time, tokens and money\\
Out of scope & other languages and benchmarks; finding the faulty code; judging whether a plausible patch is correct; tests run in parallel on many machines\\
\bottomrule
\end{tabularx}
\end{table}

\section{Structure}

The rest of the proposal is organised as follows.
\begin{itemize}[leftmargin=*,itemsep=1pt]
\item Chapter~\ref{ch:background} gives the background: the cost of a repair
session and how to order tests.
\item Chapter~\ref{ch:related} reviews related work and the open problems.
\item Chapter~\ref{ch:system} presents the validation layer and its three
modules.
\item Chapter~\ref{ch:experiments} plans the evaluation.
\item Chapter~\ref{ch:outcomes} states how each possible result will be read.
\item Chapter~\ref{ch:plan} gives the work plan, risks, threats to validity
and future work.
\end{itemize}
